\documentclass[11pt]{article}
\usepackage[margin=2.5cm]{geometry}
\usepackage{amsmath,amssymb}
\usepackage{booktabs}
\usepackage{hyperref}
\usepackage{listings}
\usepackage{xcolor}

\lstset{
  basicstyle=\ttfamily\small,
  columns=fullflexible,
  breaklines=true,
  frame=single,
  backgroundcolor=\color{gray!8}
}

\title{Optimal Strategies for the WordHunt Oracle Game\\
\large A research logbook (Volume 2)}
\author{}
\date{Started 2026-09-18}

\begin{document}
\maketitle

\begin{center}
\fbox{\begin{minipage}{0.85\textwidth}
This document is self-contained: it states the game, the model, and the
objective from scratch, and assumes no prior context. It supersedes an
earlier volume (\texttt{archive/logbook\_v1\_getting\_started.tex}) which
recorded the process of arriving at a correct model and is kept only as a
historical record -- nothing in it should be taken as current. Every
claim and number here has been independently verified against brute-force
enumeration or an independently-coded simulation, as described alongside
it; none is carried over unverified from the earlier volume.
\end{minipage}}
\end{center}

\section{The game}
\label{sec:game}

A grid of $n \times n$ cells is filled with distinct letters from an
alphabet of $n^2+1$ symbols, so that exactly one letter of the alphabet
never appears anywhere in the grid. (The real game is $n=5$: a 26-letter
alphabet, one letter absent, a $5\times5$ grid.)

A \textbf{line} is any of the $n$ rows, $n$ columns, or 2 main diagonals
-- $2n+2$ lines in total, each of length $n$. The grid is constructed so
that reading the letters along \emph{exactly one} line spells a word from
a fixed vocabulary; every other line's letters do not spell a vocabulary
word. This one-and-only-one-word property is a defining guarantee of the
real game, not an assumption we are choosing to make -- a grid violating
it is never published.

The player is told, before play starts, which single letter is absent
from the alphabet. Play then proceeds in turns: each turn, the player
names an untested letter, and the oracle discloses \emph{the one cell
where that letter actually sits in the grid} -- regardless of whether
that letter turns out to be part of the answer word or not. (The grid is
a bijection between the $n^2$ occupied cells and the $n^2$ non-absent
letters, so every query has a definite, single-cell answer; there is no
"that letter isn't part of the word" response with no location attached.)

The game ends the instant every cell of the answer line has been
individually disclosed this way -- \emph{not} at the moment the player
could, in principle, deduce the answer from partial information. If the
answer is logically forced with two of its letters still untested, the
player must still spend two more turns naming exactly those two letters
(each of which will, of course, turn out to be a hit) before the oracle
recognises the line is complete and the game reports how many turns it
took. This is the actual scoring rule of the real game, and it is the
central fact this document's objective (Section~\ref{sec:objective}) is
built around.

\section{The generative model}
\label{sec:model}

To reason about optimal play we need an explicit probabilistic model of
how a grid satisfying Section~\ref{sec:game}'s guarantee comes to exist.
We use the simplest one consistent with that guarantee, and call it
$M_2$ (the label is kept from earlier internal work so this document can
be extended later without renumbering; no other model appears in this
document).

$M_2$ generates a grid as:
\begin{enumerate}
  \item word $w \sim \mathrm{Uniform}(\text{vocabulary } W)$
  \item line $p \sim \mathrm{Uniform}(\text{the } 2n+2 \text{ lines})$
  \item excluded letter $e \sim \mathrm{Uniform}(\text{alphabet} \setminus w)$
  \item the remaining $n^2-n$ letters are placed in the remaining
  $n^2-n$ cells via a uniformly random permutation
  \item \textbf{reject and resample from step 1} if any line other than
  $p$ also happens to spell a vocabulary word
\end{enumerate}
Step 5 is what makes $M_2$ consistent with Section~\ref{sec:game}'s
guarantee: without it, a grid can accidentally spell a second word by
coincidence, and it is not merely rare -- Section~\ref{sec:n2} measures
this directly and finds it common enough, for realistic vocabulary
sizes, to matter.

The excluded letter $e$ is revealed to the player before play starts (as
in the real game, Section~\ref{sec:game}), so it is fixed per game
instance and never treated as a live hypothesis in what follows.

\textbf{Scope.} $M_2$'s vocabulary $W$ is a small, fixed, hand-picked set
of length-$n$ strings for the board sizes studied here ($n=2,3$) --
deliberately not English, just a toy dictionary large enough to be
non-trivial. Real English vocabularies (e.g.\ the $n=5$ game's
\texttt{dict\_game.csv}, 10,176 words)
are discussed only as a scale reference in Section~\ref{sec:tractability};
no result in this document has yet been computed against one under a
verified solver.

\section{The objective}
\label{sec:objective}

Let $T$ be the number of turns played until the game ends, in the sense
of Section~\ref{sec:game}: the smallest $t$ such that every cell of the
true answer line has been disclosed by one of the first $t$ queries. The
objective is to choose, adaptively (informed by every prior query and its
outcome), the sequence of letters to query so as to minimise
$\mathbb{E}[T]$ under $M_2$.

\textbf{A useful related quantity: certainty.} Separately from $T$,
define $C \le T$ as the turn at which the surviving hypothesis space --
every $(w,p)$ pair still consistent with everything disclosed so far --
first collapses to a single pair. Under $M_2$'s guarantee (unlike a model
without it), this is always well-defined and unambiguous: two distinct
$(w,p)$ pairs can never remain simultaneously consistent with a
fully-revealed grid, because the grid cannot spell two words. $C$ is not
what the game scores, and minimising $\mathbb{E}[C]$ is a different
optimisation problem from minimising $\mathbb{E}[T]$ -- Section~\ref{sec:n2}
shows the two optimal policies genuinely disagree, not just at their
optimal values but at individual decision points throughout the game. $C$
is retained here because it is analytically convenient, always $\le T$,
and useful as a baseline for measuring how much $T$'s "mandatory mop-up"
after certainty actually costs.

\section{Exact solution method}
\label{sec:method}

\textbf{State.} The relevant state for computing $\mathbb{E}[T]$ optimally
is the pair $(H, R)$: $H$, the set of $(w,p)$ hypotheses still consistent
with everything disclosed so far, and $R$, the set of grid cells disclosed
so far. Both are needed -- $R$ is not recoverable from $H$ alone, because
two different query histories can leave behind the same surviving
hypothesis set while having disclosed different amounts of whichever line
turns out to be the true one, and that difference changes $T$ directly. A
query landing on a true-line cell does double duty: it can narrow $H$
\emph{and} it always advances the mandatory reveal: a query landing on a
dross cell only narrows $H$ (if it discriminates between hypotheses at
all).

\textbf{Recursion.} Write $A \subseteq H$ for the hypotheses whose own
line is not yet fully covered by $R$ (an hypothesis whose line \emph{is}
fully covered has already "finished" -- if it were secretly true, the
game would already have ended):
\begin{align}
V(H,R) &= 0 && \text{if } A = \varnothing \\
V(H,R) &= \frac{|A|}{|H|}\Big[\,1 + \min_{L} \sum_{c} \frac{|A_c|}{|A|}\, V(A_c,\, R\cup\{c\})\,\Big] && \text{otherwise}
\end{align}
where $A_c \subseteq A$ is the subset of active hypotheses in which
letter $L$ actually sits at cell $c$, and the minimisation is over
untested letters. $\mathbb{E}[T] = V(H_0, \varnothing)$ for the initial
hypothesis set $H_0$. A letter is never excluded from the search merely
for failing to discriminate between hypotheses in $A$: it can still be
valuable purely by advancing the mandatory reveal for some or all of
them, and this objective must credit that. (A letter whose cell is
\emph{already} in $R$ must be excluded, though -- not merely discouraged
-- since considering it recurses on the identical state before that state
is cached, which is an infinite loop, not just a wasted turn.)

An analogous, simpler recursion computes $\mathbb{E}[C]$: terminate when
$H$ collapses to one $(w,p)$ pair rather than checking per-hypothesis
line coverage, with no need to track $R$ as separate state.

\textbf{Implementation.} \texttt{solver.py}: \texttt{MopUpSolver} computes
$V(H,R)$ (the real objective); \texttt{TagIdentifySolver} computes the
certainty analogue. Both are memoised expectimax searches over
exhaustively enumerated worlds (\texttt{enumerate\_worlds}, with $M_2$'s
step 5 veto implemented as a direct rejection filter on the enumerated
set). Every numeric claim below was independently cross-checked by at
least one of: an independently-coded re-derivation of the same recursion
(\texttt{simulate\_mopup\_policy\_exact\_expectation}), or literal
step-by-step play against every enumerated world as ground truth with no
expectation formula at all (\texttt{simulate\_mopup\_bruteforce}) --
\texttt{tests/test\_mopup\_solver.py} runs both at $n=2$ for every result
in Section~\ref{sec:n2} and passes.

\textbf{Greedy comparator.} A natural one-step heuristic -- at each turn,
query the letter that maximises the Shannon entropy of the immediate
outcome partition, ignoring everything beyond that one query -- is
implemented as \texttt{greedy\_choice}. It is a proxy, not a claim of
optimality: Section~\ref{sec:n2} measures exactly how much it costs to
follow it instead of the true optimum.

\section{Results at $n=2$}
\label{sec:n2}

Primary vocabulary (alphabet $\{a,b,c,d,e\}$): $W = \{\mathrm{ab, bc, cd,
de, ea}\}$, a 5-word cycle.

\subsection{How often does $M_2$'s veto actually trigger?}

Measured directly, as the fraction of $(word, line, excluded)$
combinations for which \emph{no} arrangement of the remaining letters
avoids a second word appearing (i.e.\ step 5 above always rejects):

\begin{center}
\begin{tabular}{lcc}
\toprule
vocabulary & words & impossible combinations \\
\midrule
sparse (\{ab, cd, be\}, minimal overlap) & 3 & 7.4\% \\
primary (\{ab, bc, cd, de, ea\}, a 5-cycle) & 5 & 44.0\% \\
denser (\{ab, cd, ae, bc\}) & 4 & 36.1\% \\
degenerate (every possible string counts as "a word") & 20 & 100.0\% \\
\bottomrule
\end{tabular}
\end{center}

The 100\% figure is not a bug: if literally every string is a vocabulary
word, every line of every grid trivially "spells a word", so no grid can
ever satisfy the one-word guarantee -- $M_2$ is unsatisfiable by
construction for that vocabulary, and is correctly reported as such
rather than crashing or returning something meaningless. The rate is
driven by vocabulary density (how much letter-overlap structure creates
opportunities for an accidental second word), not by anything specific to
$n=2$'s geometry -- confirmed by measurement across four vocabularies
spanning the full range from near-zero to total impossibility, rather
than assumed from the single primary case.

\subsection{Optimal values}

For the primary vocabulary, at every excluded-letter scenario (all 5 give
the same values by this vocabulary's cyclic symmetry):

\begin{center}
\begin{tabular}{ccc}
\toprule
quantity & optimal value & greedy value \\
\midrule
$\mathbb{E}[C]$ (certainty) & 2.0909 & 2.0909 \\
$\mathbb{E}[T]$ (real objective) & 2.2727 & \textbf{2.5455--3.4545} \\
\bottomrule
\end{tabular}
\end{center}

Two things worth separating out:
\begin{itemize}
  \item \textbf{The mandatory mop-up tax}: $\mathbb{E}[T] - \mathbb{E}[C]
  = 0.1818$ queries under optimal play, exactly, in every scenario for
  this vocabulary -- the measurable cost of the real scoring rule over
  merely reaching certainty.
  \item \textbf{Greedy is optimal for $C$ but not for $T$, on this
  vocabulary.} One-step maximum-entropy selection happens to match the
  true optimum for the certainty objective (2.0909 in both columns), but
  is measurably worse for the real objective -- greedy's value ranges up
  to 3.4545 depending on excluded letter, against the true optimum of
  2.2727 throughout. Greedy does not see the double-duty value of a query
  landing on the true line, and a policy tuned only for reaching
  certainty fastest can leave real mop-up queries on the table.
\end{itemize}

\subsection{The optimal policies for $C$ and $T$ genuinely disagree}

Method: solve for $\mathbb{E}[T]$ exhaustively (the solver's own search
tries every letter at every reachable state, so its memoised results
cover the whole reachable decision tree, not just the states visited
along its own optimal path). At every state where \emph{both} the
$T$-optimal and $C$-optimal solvers still faced a genuine choice (states
where the $C$-optimal solver had already reached certainty, and so had
nothing left to decide, are not a meaningful comparison -- that is
exactly the point of the mop-up tax above), compare their chosen letters,
and where they differ, compute the exact $T$-cost of forcing the
$C$-optimal choice instead, to separate real disagreement from harmless
ties (multiple equally-good letters).

Exhaustive across all 5 excluded-letter scenarios: 135 states where both
solvers faced a real choice; \textbf{65 disagreed (48\%), and all 65 were
confirmed genuine} -- forcing the $C$-optimal choice measurably costs
more $T$ in every one, none are ties. Worst single-state gap: exactly 1.0
extra expected query.

\textbf{Conclusion: these are different policies, not just different
values.} A solver built to reach certainty fastest is not a substitute
for one built to minimise the real scoring rule, at almost half of all
decision points in this small a game.

\section{Results at $n=3$}
\label{sec:n3}

Vocabulary (alphabet $\{a,\ldots,j\}$): $W = \{\mathrm{abc, cde, efg,
ghi, ijb, bad, dec, fah}\}$.

\subsection{Optimal values, all ten excluded-letter scenarios}

\begin{center}
\begin{tabular}{ccccc}
\toprule
excluded & $M_2$ hypotheses & $\mathbb{E}[C]$ & $\mathbb{E}[T]$ & mop-up tax \\
\midrule
a & 27{,}420 & 3.3700 & 4.0161 & 0.6461 \\
b & 27{,}558 & 3.3060 & 3.9792 & 0.6733 \\
c & 26{,}922 & 3.5938 & 4.2672 & 0.6734 \\
d & 27{,}462 & 3.5741 & 4.2423 & 0.6682 \\
e & 27{,}144 & 3.2555 & 3.8767 & 0.6212 \\
f & 32{,}376 & 3.5607 & 4.1069 & 0.5463 \\
g & 32{,}126 & 3.5576 & 4.1672 & 0.6096 \\
h & 32{,}424 & 3.4869 & 4.1172 & 0.6302 \\
i & 32{,}100 & 3.4886 & 4.1031 & 0.6145 \\
j & 37{,}014 & 3.8287 & 4.4244 & 0.5956 \\
\bottomrule
\end{tabular}
\end{center}

The mop-up tax at $n=3$ (0.55--0.67 queries) is roughly $3$--$4\times$ the
$n=2$ tax (0.18) -- two data points, not yet a trend to extrapolate from,
but consistent with "more cells per line to mop up costs more."

\subsection{Policy divergence, sampled}

Repeating Section~\ref{sec:n2}'s comparison exhaustively at $n=3$ is not
practical (millions of memoised states per scenario); a bounded random
sample of 3000 states for excluded=`a' found 142 states where both
solvers faced a genuine choice, of which \textbf{138 (97\%) were
confirmed genuine divergences}, none ties, worst gap again exactly 1.0
extra expected query. Consistent with, and if anything stronger than, the
$n=2$ result.

\subsection{Greedy vs.\ optimal}

\begin{center}
\begin{tabular}{ccccc}
\toprule
excluded & $C$-optimal & $C$-greedy & $T$-optimal & $T$-greedy \\
\midrule
a & 3.3700 & 5.1041 & 4.0161 & 8.3436 \\
b & 3.3060 & 5.1876 & 3.9792 & 8.3441 \\
c & 3.5938 & 5.3325 & 4.2672 & 8.3834 \\
d & 3.5741 & 5.3069 & 4.2423 & 8.3925 \\
e & 3.2555 & 5.1746 & 3.8767 & 8.3531 \\
f & 3.5607 & 5.3824 & 4.1069 & 8.3001 \\
g & 3.5576 & 5.3590 & 4.1672 & 8.2653 \\
h & 3.4869 & 5.4239 & 4.1172 & 8.3020 \\
i & 3.4886 & 5.4283 & 4.1031 & 8.3271 \\
j & 3.8287 & 5.5378 & 4.4244 & 8.3437 \\
\bottomrule
\end{tabular}
\end{center}

Unlike at $n=2$, greedy is now measurably worse than optimal for
\emph{both} objectives, not just $T$: roughly $+1.7$ queries for
certainty, and consistently landing near $8.3$ against a true optimum of
$3.9$--$4.4$ for the real objective -- roughly \textbf{double} the
optimal value, dramatically larger than the $n=2$ gap (2.27 vs.\ up to
3.45). $n=2$'s "greedy matches optimal for certainty" coincidence
(Section~\ref{sec:n2}) does not survive to $n=3$: with more structure to
exploit, one-step maximum-entropy selection is a substantially worse
proxy for both objectives studied here, not a merely-occasionally-worse
one.

\section{Where exact solving stops being tractable}
\label{sec:tractability}

Two distinct obstacles, easy to conflate:

\textbf{1. Full-world enumeration has a factorial wall.} Computing
$V(H,R)$ exactly as defined requires the hypothesis space $H$ to range
over literal, fully-populated grids (so that every query's outcome cell
is well-defined for every hypothesis). The number of ways to permute the
$n^2-n$ dross letters into the remaining cells is $(n^2-n-1)!$:

\begin{center}
\begin{tabular}{cccc}
\toprule
$n$ & alphabet size & lines & dross permutations $(n^2-n-1)!$ \\
\midrule
2 & 5  & 6  & $2! = 2$ \\
3 & 10 & 8  & $6! = 720$ \\
4 & 17 & 10 & $12! = 479{,}001{,}600$ \\
5 & 26 & 12 & $20! \approx 2.43\times10^{18}$ \\
\bottomrule
\end{tabular}
\end{center}

$n=4$ is already far beyond reach for literal enumeration. \textbf{No
valid way around this is currently implemented.} A natural idea --
represent a hypothesis by its $(w,p)$ pair alone, without the dross
permutation, since the permutation's count is the same for every
hypothesis and should cancel out of relative probabilities -- was tried
and found \emph{not} to work for this objective: a dross letter's actual
disclosed cell depends on which specific permutation is realised, which
is not determined by $(w,p)$ alone, and the distribution over which cell
it could be stops being uniform across hypotheses once some cells are
already disclosed. A correct escape from full-world enumeration would
need to properly count, for a given $(w,p)$ and disclosure history $R$,
how many consistent completions place a given untested letter at each
possible cell -- a real combinatorial counting problem, not yet solved.

\textbf{2. Even where full-world enumeration is possible, the search
itself has a large branching factor.} Every query resolves to one of up
to $n^2$ distinct cells (not a small handful), so the reachable-state
count under exhaustive expectimax search grows fast even at $n=3$
(3.4--4.5 million memoised states per scenario for \texttt{MopUpSolver}
in Section~\ref{sec:n3}, taking 45--80 seconds each on ordinary
hardware). This is a separate obstacle from (1): it doesn't go away even
once $H$ is confined to small, tractable vocabularies, and it will only
get worse at $n=4$ (already-infeasible for reason (1)) and dramatically
worse at $n=5$ (up to 21 distinct queryable letters, and a real
dictionary's vocabulary is $\sim$10,000 words rather than a handful).

\textbf{$n=5$ at real-dictionary scale has not yet been validly
attempted.} Given obstacle (1) already blocks full enumeration at $n=4$,
and no correct escape from it exists yet, there is currently no way to
even \emph{start} an exact computation at $n=5$'s real scale, independent
of obstacle (2)'s branching-factor cost. This is stated as an open gap,
not a measured wall -- an earlier attempt at measuring this wall used a
solver later found to be invalid for this objective, and no valid
replacement measurement has been made.

\section{Open questions}
\label{sec:next}

\begin{enumerate}
  \item \textbf{A correct pair-based (non-full-world) solver}, per
  Section~\ref{sec:tractability}(1) -- the highest-value next step, since
  it is the only route past $n=4$'s factorial wall at all.
  \item \textbf{Branch-and-bound with a provable lower bound}, to reduce
  the cost from Section~\ref{sec:tractability}(2) without sacrificing
  exactness.
  \item \textbf{Solving from mid-game states rather than the opening.} A
  real player starts from wherever they currently are, typically with
  $|H|$ already reduced by prior disclosures; the opening state is the
  single largest and hardest one the game ever presents. Whether
  mid-game states are tractable further into $n=4$ or $n=5$ than the
  opening is, is worth checking directly rather than assuming either way.
  \item \textbf{A variance-sensitive objective}, e.g.\ minimising
  $\mathbb{E}[T^2]$ rather than $\mathbb{E}[T]$ -- requires tracking two
  moments through the recursion instead of one; not yet attempted.
\end{enumerate}

\end{document}
